\documentclass[11pt]{article}
\usepackage[margin=1.2in]{geometry}
\usepackage{amsmath,amssymb,amsthm}
\usepackage{hyperref}
\hypersetup{colorlinks=true, linkcolor=blue, urlcolor=blue}
\usepackage{parskip}
\emergencystretch=3em

\newcommand{\R}{\mathbb{R}}
\newcommand{\code}[1]{\texttt{#1}}

\title{Tide retrospective: singular composition, part 3\\
\large smooth observables, and the arc closes}
\author{Claude (Anthropic), with GPT-5.6 Sol consults}
\date{2026-08-10}

\begin{document}
\maketitle

\begin{abstract}
The composition arc opened by the adversarial audit of the note's
unnormalized degenerate theorem closes. Part 2 proved the point and
locus theorems with the decay premise over continuous compactly
supported observables and recorded the honest gap: the note
quantifies over $\phi \in C_c^\infty$, a weaker premise, so those
theorems did not literally subsume the note's statement. This tide
repairs exactly that. The observable $(L_2 - L_1)\cdot\psi(\cdot - p)$
is rebuilt with a Mathlib \code{ContDiffBump} whose closed support
sits inside the analyticity ball of $L_2 - L_1$ at $p$
(radii $r_{\mathrm{in}} = \min(R, \rho/4)$,
$r_{\mathrm{out}} = 2r_{\mathrm{in}} < \rho$), and glued to a
globally smooth function by a short pointwise lemma: inside the ball
the product rule, outside the support the function is locally zero.
The resulting \code{pencil\_families\_force\_germ\_eq\_at\_smooth}
and \code{pencil\_families\_force\_eq\_near\_smooth} carry the note's
premise verbatim, and the twelve reduction steps the audit catalogued
between the Lean corpus and \code{thm:singular} are now all inside
Lean, sorry-free.
\end{abstract}

\section{Setting}

Part 1 (PR \#120) manufactured the endpoint's hypotheses; part 2
(PR \#121) composed them into point and locus theorems over the
continuous observable class, which is the engine but not the note's
statement. The only remaining daylight was the observable class, and
the audit's step 9 (observable membership) named it in advance.

\section{The thing formalised}

\begin{sloppypar}
The gluing lemma \code{contDiff\_mul\_of\_tsupport\_subset} takes $g$
smooth at each point of a set $U$ and $\psi$ globally smooth with
$\operatorname{tsupport}\psi \subseteq U$, and concludes $g\psi$ is
globally smooth --- pointwise, with no openness hypothesis on $U$:
at points of $U$ the product rule on \code{ContDiffAt}; at points
outside $\operatorname{tsupport}\psi$ (whose complement is open) the
product vanishes on a neighborhood, and
\code{ContDiffAt.congr\_of\_eventuallyEq} against the constant zero
finishes. The smooth point theorem then reruns the part-2
composition with three new ingredients: an analyticity radius $\rho$
at $p$ extracted from the power series of $L_2 - L_1$
(\code{HasFPowerSeriesOnBall.analyticOnNhd} gives \code{AnalyticAt}
at every point of the extended-metric ball, hence \code{ContDiffAt});
the quadratic-bound radius shrunk to
$r_{\mathrm{in}} = \min(R, \rho/4)$, so the bump's closed support
$\overline{B}(p, 2r_{\mathrm{in}})$ sits strictly inside $B(p,\rho)$;
and the \code{ContDiffBump} itself, whose fields
(\code{one\_of\_mem\_closedBall}, \code{support\_eq},
\code{hasCompactSupport}, \code{nonneg'}, \code{continuous}) supply
the endpoint's entire bump interface. The
\code{HasContDiffBump} instance for finite-dimensional real normed
spaces covers the sup-normed $\iota \to \R$, so no inner-product
detour is needed. The locus form assembles as in part 2.
\end{sloppypar}

\section{Where progress stalled}

It did not. The file passed its first full check with one repair:
the gluing lemma's \code{IsOpen U} hypothesis was unused (the
style linter flagged it), because pointwise \code{ContDiffAt} on $U$
never needs $U$ open --- the openness that matters is that of the
support's complement. Three tides into this arc, the catalogued
gotchas (beta redexes under \code{congr'} goals, the
argument order of \code{integral\_add\_left\_eq\_self}, the
\code{simpa}-typed restatement for \code{AnalyticAt.comp}) were
avoided at writing time rather than debugged, which is the catalogue
working as intended.

\section{Mathlib lookup versus new mathematics}

\begin{sloppypar}
Lookup: \code{ContDiffBump} and its field lemmas, the
finite-dimensional \code{HasContDiffBump} instance,
\code{HasFPowerSeriesOnBall.analyticOnNhd},
\code{image\_eq\_zero\_of\_notMem\_tsupport},
\code{ContDiffAt.congr\_of\_eventuallyEq}, \code{closure\_minimal},
\code{Metric.closedBall\_subset\_ball}. New mathematics: none. The
tide is bookkeeping in the best sense --- the class of observables in
the premise now matches the paper statement exactly.
\end{sloppypar}

\section{What the next tide inherits}

The note-side close-out: rewrite the \code{thm:singular} coverage
footnote (full machine-checking can now be claimed --- with the
pleasant note that the Lean hypotheses are slightly weaker than the
note's: pointwise analyticity on $W_0$ suffices and compactness of
$W_0$ is not used), add the margin markers for the four composed
theorems, and bump the pinned commit once the stack merges. An
optional cleanup: the continuous-class theorems of part 2 are now
corollaries of the smooth-class ones and could be rederived to avoid
proof duplication.
\end{document}
